%! TeX root = ./main.tex
% vim: set spelllang=fr:
\selectlanguage{french}
\section{Comparaison et perspectives}

% TODO: introduire la question de recherche ici

Dans ce chapitre ont été présentées deux approches généralisant la réconciliation bilatérale à des contextes où plusieurs unités dépendantes évoluent conjointement.
Ce cas de figure se présente lorsque de multiples gènes évoluent ensemble le long d’un génome.
Le principe général consiste à intégrer un niveau intermédiaire dans la hiérarchie évolutive, décrivant l’évolution d’adjacences (DeCo) ou de segments porteurs d’ensembles d’unités (super-réconciliation).

L’approche DeCo permet de modéliser des événements qui affectent des paires d’unités voisines, incluant notamment des réarrangements.
D’un point de vue algorithmique, elle consiste à réconcilier individuellement les arbres phylogénétiques de chaque unité, puis à rassembler des paires de nœuds pour former une forêt d’adjacences.
Chaque racine de cette forêt correspond au gain d’une adjacence qui se propage ensuite dans sa descendance.
L’objectif est de minimiser le nombre de racines, c’est-à-dire le nombre de gains d’adjacences.
Le statut de complexité du problème est encore inconnu, quoiqu’il ait été conjecturé NP-complet, mais un algorithme polynomial existe lorsque les réconciliations individuelles sont fixées en avance.
La modélisation à l’échelle des adjacences impose des limites sur les événements pouvant être pris en compte, puisqu’elle exclut les événements qui affectent simultanément plus de deux unités.
Par ailleurs, un résultat optimal peut comporter des structures non-linéaires, dans lesquelles une unité est voisine de trois autres ou plus, incompatibles avec les structures linéaires ou circulaires des génomes.

La super-réconciliation propose un paradigme similaire, basé sur l’évolution de segments joignant un nombre quelconque d’unités.
Cette extension permet de modéliser des événements segmentaux affectant plusieurs unités.
De la même façon que pour DeCo, l’objectif algorithmique est de construire un arbre dans lequel chaque nœud rassemble des segments de nœuds issus de réconciliations individuelles.
La structure reconstruite ici est un arbre et non une forêt car tous les segments sont supposés dériver d’un segment ancestral unique.
Le problème d’optimisation combinatoire associé est NP-complet dans sa version générale, lorsque la topologie de l’arbre de synténies est inconnue.
Si la topologie est fixée, la complexité est inconnue dans la déclinaison du problème avec des segments, et polynomiale dans sa déclinaison avec des groupes.
Les modèles d’évolution considérés jusqu’à présent pour la super-réconciliation sont encore limités.
Contrairement à ceux considérés sous DeCo, ils ne permettent par exemple pas de modéliser des duplications en tandem, des fusions ou fissions de synténies ou des réarrangements.

\subsection{Thèse}

L’objectif général du projet de recherche est de développer la théorie de la réconciliation multiple, en particulier autour de l’approche de super-réconciliation, et de démontrer son potentiel dans l’analyse de l’évolution de systèmes biologiques.

% Contributions
